\documentclass[11pt,letterpaper]{article}

% Academic CV for Fall 2027 PhD outreach
% Candidate: Kazi Md. Tawsif Rahman

\usepackage[letterpaper,margin=0.68in]{geometry}
\usepackage[T1]{fontenc}
\usepackage[utf8]{inputenc}
\usepackage{lmodern}
\usepackage{enumitem}
\usepackage[hidelinks]{hyperref}
\usepackage{titlesec}
\usepackage{xcolor}
\usepackage{tabularx}
\usepackage{array}
\usepackage{microtype}
\usepackage{fancyhdr}

\input{glyphtounicode}
\pdfgentounicode=1

\definecolor{darkblue}{HTML}{17365D}

\hypersetup{
  pdftitle={Academic CV -- Kazi Md. Tawsif Rahman},
  pdfauthor={Kazi Md. Tawsif Rahman},
  pdfsubject={Academic curriculum vitae},
  pdfkeywords={AI security, network security, privacy-preserving machine learning, membership inference, differential privacy, trustworthy machine learning}
}

\titleformat{\section}
  {\large\bfseries\color{darkblue}}
  {}
  {0em}
  {}
  [\titlerule]

\titlespacing*{\section}{0pt}{8pt}{3pt}

\setlength{\parindent}{0pt}
\setlength{\parskip}{2pt}
\setlength{\tabcolsep}{0pt}
\setlist[itemize]{
  leftmargin=1.3em,
  label=\textbullet,
  itemsep=1.8pt,
  parsep=0pt,
  topsep=2pt
}
\setlist[enumerate]{
  leftmargin=1.55em,
  itemsep=2.5pt,
  parsep=0pt,
  topsep=2pt
}

\pagestyle{fancy}
\fancyhf{}
\renewcommand{\headrulewidth}{0pt}
\renewcommand{\footrulewidth}{0pt}
\fancyfoot[C]{\footnotesize Kazi Md. Tawsif Rahman \textbar\ Academic CV \textbar\ \thepage}

\newcommand{\entry}[2]{%
  \begin{tabularx}{\textwidth}{@{}X >{\raggedleft\arraybackslash}p{0.19\textwidth}@{}}
    \textbf{#1} & #2
  \end{tabularx}\par
}

\newcommand{\projectentry}[3]{%
  \textbf{#1}\par
  \textit{#3}\hfill\textit{#2}\par
}

\begin{document}

{\LARGE\bfseries Kazi Md. Tawsif Rahman}\\[2pt]
Dhaka, Bangladesh \enspace\textbar\enspace
\href{mailto:tawsifcse113@gmail.com}{tawsifcse113@gmail.com} \enspace\textbar\enspace
+880 1533-406937 \enspace\textbar\enspace
\href{https://linkedin.com/in/tawsifrahman113}{linkedin.com/in/tawsifrahman113}\\
\href{https://github.com/tawsif113}{github.com/tawsif113} \enspace\textbar\enspace
\textbf{Research portfolio:}
\href{https://research.tawsifrahman.flaro-tech.com/}{research.tawsifrahman.flaro-tech.com}

\section*{Research Summary}

Computer Science graduate and backend systems engineer researching differential
privacy and membership-inference risk in machine-learning-based network intrusion
detection. My current independent study uses Opacus-based DP-SGD, explicit
$(\varepsilon,\delta)$-privacy accounting, and shadow-model membership-inference
auditing under a locked evaluation protocol. I completed the five-seed NSL-KDD
analysis and a constrained single-seed UNSW-NB15 external validation. The
experimental roadmap is frozen and a venue-neutral first paper draft is under
author review.

\section*{Research Interests}

AI security; differential privacy; privacy-preserving machine learning;
membership-inference auditing; trustworthy evaluation of security-critical ML;
federated learning for security applications; reproducible experimental systems.

\section*{Education}

\entry{Chittagong University of Engineering and Technology (CUET), Bangladesh}{July 2024}
B.Sc. in Computer Science and Engineering \hfill CGPA: 3.59 / 4.00

\section*{Research Experience}

\projectentry{Privacy--Utility Auditing of DP-SGD for ML-Based Network Intrusion Detection}
{2026 -- Present}{Independent Research Project}
\begin{itemize}
  \item Designed a locked NSL-KDD binary-classification protocol with a
  70/10/20 target-train/validation/shadow split, validation-only F2 threshold
  selection, and untouched KDDTest+ utility evaluation; RF, XGBoost, and MLP
  utility baselines are retained as archived context.

  \item Built five shadow MLPs and evaluated score-only and label-aware
  membership-inference attacks. The strongest evaluated non-private attack
  remained near chance, creating a floor effect that is reported as an
  evaluation limitation rather than evidence of privacy.

  \item Reimplemented the target MLP in PyTorch and integrated Opacus DP-SGD
  with explicit privacy accounting, clipping, sampling, noise, and training
  parameters recorded for reproducibility.

  \item Completed five target-training seeds for non-private,
  $\varepsilon\approx4$, and $\varepsilon\approx2$ on NSL-KDD. Mean Recall was
  70.93\%, 71.34\%, and 70.11\%, respectively; the $\varepsilon\approx4$
  Recall difference was uncertain, while false-positive rate increased and
  average precision decreased.

  \item Across the accepted NSL-KDD analysis, paired membership-inference AUC
  intervals crossed zero, so the evidence does not support a claim that DP-SGD
  reduced measured overall membership leakage.

  \item Completed a constrained single-seed UNSW-NB15 validation. The private
  settings preserved selected-threshold F1 and Recall closely but increased the
  false-positive burden and reduced average precision; overall MIA remained near
  chance.

  \item The experimental roadmap is complete. A venue-neutral first manuscript
  is frozen for author review and target-venue selection; no tested epsilon is
  claimed to be universally optimal.
\end{itemize}
\textit{Code, results, and reproducibility artifacts:}
\href{https://github.com/tawsif113/privacy-utility-dp-ids}{github.com/tawsif113/privacy-utility-dp-ids}

\section*{Publications}

\begin{enumerate}
  \item \textbf{Kazi Md. Tawsif Rahman} and Chowdhury Mahfuzul Hoq.
  ``An Automated System for Detecting Property Insurance Fraud Using Machine
  Learning.'' \textit{International Journal of Mathematical Sciences and
  Computing}, 2024.
  \href{https://doi.org/10.5815/ijmsc.2024.03.02}{doi:10.5815/ijmsc.2024.03.02}.
  \par\smallskip

  \item Mahfuzulhoq Chowdhury, \textbf{Kazi Md. Tawsif Rahman}, and Hossain
  Ahmad Maruf. ``Whistle Blower: An Insurance Awareness Mobile Application with
  Insurance Policy Selection, Fraud Detection, Critical Help, Complaint
  Features.'' \textit{2024 8th International Conference on Computational System
  and Information Technology for Sustainable Solutions (CSITSS)}, pp. 1--6,
  2024.
  \href{https://doi.org/10.1109/CSITSS64042.2024.10817002}{doi:10.1109/CSITSS64042.2024.10817002}.
\end{enumerate}

\section*{Selected Professional Experience}

\entry{Software Engineer, BRAC IT Services Ltd., Dhaka, Bangladesh}{Dec 2024 -- Present}
\textbf{CRM Platform}
\begin{itemize}
  \item Built configurable rule-based routing with priorities, schedules,
  conditional matching, rotation rules, duplicate checks, and validated
  field/operator/value conditions.

  \item Replaced full-column Kanban reordering with fixed-length order keys and
  reduced each movement to a single-rank assignment.

  \item Added access controls, idempotent moves, transition validation, audit
  logs, and timeline entries to make state changes traceable.
\end{itemize}

\textbf{Ticketing System}
\begin{itemize}
  \item Built services for ticket creation, search, comments, forwarding,
  assignment rules, teams, SLA/KPI tracking, feedback, notifications, and
  metadata.

  \item Implemented least-loaded team routing with Redis-backed workload
  counters and maintained PostgreSQL--Redis workload consistency after ticket
  creation, reassignment, and resolution.
\end{itemize}

\section*{Technical and Research Skills}

\begin{tabularx}{\textwidth}{@{}>{\bfseries}p{0.20\textwidth} X@{}}
Research methods &
Locked-split experimental design, validation-only threshold selection,
shadow-model membership-inference auditing, privacy accounting, paired bootstrap
confidence intervals, reproducibility manifests, threat-model specification \\[1pt]
ML / Privacy &
Python, scikit-learn, XGBoost, PyTorch, Opacus, DP-SGD, membership-inference
attacks, model evaluation \\[1pt]
Backend / Data &
Java, Spring Boot, Spring Security, JPA/Hibernate, REST APIs, PostgreSQL,
MongoDB, Redis, RabbitMQ, Flyway \\[1pt]
Engineering &
SQL, C, C++, Docker, Linux, Git, JUnit 5, Mockito, API testing
\end{tabularx}

\section*{Awards and Achievements}

\begin{itemize}
  \item Codeforces Specialist; rating above 1500.
  \item Finalist, CUET IUPC 2024 and CUET IUPC 2023.
  \item Finalist, CUET Tech Carnival 2020.
\end{itemize}

\section*{Selected Coursework}

Algorithms and Data Structures; Database Systems; Computer Networks; Operating
Systems; Software Engineering; Artificial Intelligence; Machine Learning;
Information Security.

\end{document}
